\section{Conclusion}

In recent years, the combination of structured knowledge representation and large language models has opened new opportunities for building systems that are both intelligent and interpretable. Ontologies and knowledge graphs play a crucial role in this space, providing the foundation for organizing complex domain knowledge in a transparent and reusable way. When integrated with generative models, they enable a new form of reasoning, one that combines the precision of structured data with the flexibility of natural language understanding.
In domains such as food, nutrition, and sustainability, this integration is particularly important. These areas require not only factual accuracy, but also contextual understanding, scientific traceability, and sensitivity to public health and environmental concerns. Structuring information about ingredients, nutrients, and health outcomes in a formal schema allows knowledge to be shared, queried, and validated across systems, supporting more consistent and evidence-based communication.
More broadly, this direction illustrates a shift toward AI systems that are not just powerful, but explainable and accountable. By grounding model outputs in well-defined knowledge structures, we move closer to applications that support education, policy, and decision-making in transparent and trustworthy ways.